\section{Introduction}\label{sec:intro}

Data pipelines fail in ways that are cheap to describe and expensive to fix: an upstream schema
changes type, a source stalls, a co-tenant starves compute, ingress bursts past capacity. The
standard remedy is an on-call human, whose response time dominates recovery time. Kirubakaran et
al.~\cite{kirubakaran2025governing} propose replacing part of that loop with four bounded LLM agents
(monitoring, optimisation, schema, recovery) whose only output is a structured action proposal,
gated by a policy engine before anything executes, and report reductions of up to 45\% in recovery
time, about 25\% in cost, and over 70\% in manual interventions against static orchestration.

This paper asks whether that claim survives an evaluation in which the agents are a real, remote
language model rather than a simulation, in which the baselines include cheap automation rather than
only a slow human, and in which the safety guarantee is attacked rather than asserted. It is a
replication and extension of \cite{kirubakaran2025governing}, together with the engineering record
of taking the resulting system from a research artefact to something an operator could run.

\paragraph{Thesis.}
Governance claims hold only if the evaluation path is the production path. Every headline number in
the replicated paper's family of results depends on choices about \emph{how} the system is exercised
--- a mock or a live model, a modelled or a measured human, an assumed or a measured provisioning gap
--- and each of those choices, made for tractability, silently changes what is being measured. We
make that dependence measurable rather than arguing it away.

\paragraph{Contributions.}
\begin{enumerate}
  \item \textbf{A live-model replication} on a full matrix of eight configurations (three
    non-agent comparison points, the full system, and four single-agent ablations) over four fault
    scenarios, with per-run provenance and billing-accurate LLM accounting (\S\ref{sec:method}).
  \item \textbf{A measured mock-versus-live divergence}: the same configuration under a
    deterministic mock and under a live model, compared directly (\S\ref{sec:results}).
  \item \textbf{Credible baselines and decision quality}: rule-based and autoscaling comparators,
    and a correct-mitigation metric in addition to speed.
  \item \textbf{An adversarial evaluation with an independent specification oracle}, reporting
    containment, over-blocking and exact agreement with Wilson intervals, and the defects it found
    (\S\ref{sec:results}, \S\ref{sec:lessons}).
  \item \textbf{Closed-form sensitivity analyses} for the two modelled assumptions the headline
    numbers rest on --- the simulated human's latency and the cost model's provisioning gap.
  \item \textbf{Production-hardening lessons}, each tied to a defect measured in this system and to
    the design record (\S\ref{sec:lessons}).
\end{enumerate}

\paragraph{What is not novel.}
The architecture --- three planes, four bounded agents, a policy gate, an observe--reason--propose
loop --- is that of \cite{kirubakaran2025governing}; we implement it, we did not invent it. We claim no
new agent design, no new policy language, and no new statistical method. The contribution is the
evidence, the measurement discipline, and the defects found by holding the system to its own claims.

\paragraph{Roadmap.}
Section~\ref{sec:related} places the work; \S\ref{sec:system} describes the system as implemented;
\S\ref{sec:method} the experimental design; \S\ref{sec:results} the results;
\S\ref{sec:lessons} the hardening lessons; \S\ref{sec:discussion} limitations and threats to
validity.
